%! Unit Launch: Study Design — Unit 1, Lesson 1.0 — Lesson Plan
% PILOT SHAPE (2026-09-06): modeled on AP Statistics lesson 1.4 minus its AP Practice page.
% GRADUAL-RELEASE plan (warm-up / guided notes and practice / spoken debrief / close and START
% THE HOMEWORK). The guided notes carry the release ROW BY ROW in one Main Ideas / Notes table
% (2026-09-12 shape): I do / we do in rows 1-4 (problems 1-13) -> we do the Guided Practice row
% (problems 14-17). THERE IS NO INDEPENDENT PRACTICE SET IN THE NOTES — the homework is the
% individual practice, and the period ends with students starting it. Teacher-facing. Breakable
% boxes are `skillbox` with a \boxguard ahead; the two tabularx boxes are `fixedskillbox` (never
% splits). No group activity, no tier boxes, no exit ticket. (Converted from the group-activity
% shape 2026-09-14: the meet-card activity's crux — the bib-number claim — is the debrief's cold
% check; its Student Life Survey practice box is the Guided Practice row.)
\documentclass[10pt]{article}
\usepackage{saar-article}
\usepackage{saar-boxes}
\usepackage{graphicx}   % -article does NOT load graphicx
\graphicspath{{images/}}
\newcommand{\TallMath}[1]{$\displaystyle #1\rule[-1.4em]{0pt}{3.2em}$}
\newcommand{\UnitNumberName}{Unit 1: Asking Questions with Data: Study Design \quad}
\newcommand{\LessonNumberName}{Lesson 1.0: Unit Launch: Study Design}

\begin{document}
\begin{center}
    {\Large\bfseries \CourseName \\
    {\small\bfseries \UnitNumberName \LessonNumberName}}
\end{center}

\begin{tcolorbox}[colback=frost, colframe=cerulean, boxrule=0.9pt, arc=2mm,
    left=2mm, right=2mm, top=1.5mm, bottom=1.5mm]
    \textbf{Primary Objective:} Students will explain that data are numbers \emph{with a
    context}, identify the individuals and variables in a data set, classify each variable
    as categorical or quantitative and justify the call, apply the digit test to a column of
    digits that cannot be averaged, and recognize a statistical question --- justifying every
    choice from what the column records, never from what it looks like.
    \\[2pt]
    \textbf{Standards:} PS.DC.1a, PS.DC.1b, PS.DC.1c; AFDA.DA.2a \quad$\cdot$\quad
    \emph{previews} PS.DC.1d--e (1.2), PS.DC.2 (1.3), PS.DC.3 (1.5--1.7)
    \\[2pt]
    \textbf{Lesson model:} \emph{Gradual release --- I do, we do, you do --- all of it inside
    the guided notes.} There is no group activity. The notes name the vocabulary as they teach
    it and deliver the instruction \textbf{row by row in one Main Ideas / Notes table} (four
    rows, problems 1--13: you read the definition, mark up the display, and work the first
    problem of each grid; the class works the rest); the Guided Practice row (problems 14--17)
    is worked together with students holding the pen; \textbf{the homework is the individual
    practice}, and students start it in class while you circulate. The whole-class debrief is
    \emph{spoken} --- there is no practice set at the end of the notes and no exit ticket.
    \\[2pt]
    \textbf{Note on the unit launch:} this is Lesson 0 --- \emph{population} and \emph{sample}
    are used in plain English today (warm-up item~2, homework item~5) and named formally in
    Lesson~1.2. Do not introduce the words; let the idea sit.
\end{tcolorbox}

\renewcommand{\arraystretch}{1.4}
\boxguard
\begin{skillbox}[Learning Targets \& Key Understandings]{goldbox}
\begin{tabularx}{\linewidth}{@{}X|X@{}}
\textbf{Learning targets (student language)} & \textbf{Key understandings (the ``why'')} \\[2pt]
\hline
\vspace{-6pt}
\begin{itemize}[leftmargin=1.1em, nosep, after=\vspace{2pt}]
  \item I can explain why a number is not \textbf{data} until I know who was measured, what
        was measured, and in what units. \hfill \textit{PS.DC.1a--b}
  \item I can read a data table and name the \textbf{individuals} (rows) and the
        \textbf{variables} (columns). \hfill \textit{PS.DC.1c; AFDA.DA.2a}
  \item I can classify a variable as \textbf{categorical} or \textbf{quantitative} and
        justify it from what the column records --- an amount or a label.
        \hfill \textit{PS.DC.1c; AFDA.DA.2a}
  \item I can use the \textbf{digit test} on a column of digits, summarize each kind the way
        its type allows, and recognize a \textbf{statistical question}.
        \hfill \textit{PS.DC.1a--c}
\end{itemize}
&
\vspace{-6pt}
\begin{itemize}[leftmargin=1.1em, nosep, after=\vspace{2pt}]
  \item \emph{Data are collected for a purpose and have meaning in a context.} ``$12$''
        answers nothing; ``the six players on the roster scored a mean of $12$ points''
        answers something.
  \item The \textbf{type} of a variable decides what you are allowed to do with it. You may
        average a quantitative variable; a categorical variable can only be counted and turned
        into a percent.
  \item \textbf{Target misconception: the digit reflex.} Students classify anything written
        with digits as quantitative. Jersey number, grade level, ZIP code, and Member ID are
        categorical. The test is not ``does it look like a number?'' but ``does arithmetic on
        it mean anything?'' --- the blog's $84/6 = 14$ is correct and describes nobody.
  \item \textbf{Second misconception:} after the digit test lands, some students swing the
        other way and reclassify Height or Commute as categorical. Digits are not the enemy;
        $70$ inches is an amount.
  \item A \textbf{statistical question anticipates variability}: ``How tall is Marcus?'' has
        one answer; ``How tall are the players?'' has many.
\end{itemize}
\end{tabularx}
\end{skillbox}

\boxguard[24]
\begin{skillbox}[{Vocabulary, Concepts, \& Theorems --- taught directly in the guided notes}]{greenbox}
\small
\renewcommand{\arraystretch}{1.35}
\begin{tabularx}{\linewidth}{@{} >{\bfseries}p{3.1cm} X @{}}
Data & Numbers or labels collected about individuals, together with the context --- who,
       what, units --- that gives them meaning. \\
Individual & One person, animal, or object that data are collected about --- one \emph{row}
       of a data table. \\
Variable & A characteristic recorded for each individual --- one \emph{column} of a data
       table. The name column is not one. \\
Categorical variable & A variable that records a \emph{label} or group name (position, lunch
       choice, grade level, yes/no). Summarize with counts and percents. \\
Quantitative variable & A variable that records a \emph{measured amount} with units (points,
       minutes, siblings). Summarize with a mean or median. \\
Digit test & The check that settles a column written with digits: \emph{does arithmetic on
       this variable mean anything?} \\
Statistical question & A question whose answers are expected to \emph{vary} from individual
       to individual, so it takes data to answer. \\
Population, sample & The whole group you want to know about; the part you actually collect
       data from. \emph{(Used in plain English today; named formally in Lesson 1.2.)} \\
\end{tabularx}
\end{skillbox}

% fixedskillbox never splits, so the phase table stays intact on one page.
% THE PHASE TABLE IS A CONTRACT: 5 / 35 / 10 / 10 — the I do is ~20 of the 35 and the Guided
% Practice ~15; the debrief is spoken; the homework start is the second formative read.
\begin{fixedskillbox}[Lesson at a Glance --- \MeetingLength]{frost}
\renewcommand{\arraystretch}{1.25}
\begin{tabularx}{\linewidth}{@{}l c X X@{}}
\textbf{Phase} & \textbf{Min} & \textbf{Students} & \textbf{Teacher} \\
\hline
Warm-Up & 5 & Three Algebra 1 items --- a percent, a scaled-up survey, a mean with a
  sentence --- silent & Debrief item~3(b) aloud; write the sentence on the board and leave
  it up all period \\
Guided Notes \& Practice & 35 & Fill the vocabulary box as each term is named; work rows
  1--4 of the table (problems 1--13) with the pen in their hand; then the Guided Practice row,
  \emph{Student Life Survey} (14--17), together & \textbf{I do / we do, row by row}: read the
  definition, mark up the display, work problems 1, 3, 6, 11; the class works the rest (20)
  $\to$ \textbf{we do} Guided Practice 14--17 (15) \\
Debrief & 10 & Pens down, papers up --- answer aloud, nothing to fill in & Put the blog's
  $14$ back on the board and take the digit reflex, then the ``it's numbers'' reason; cold
  check on the meet-card bib claim \\
Close \& Assign & 10 & \textbf{Start the homework in class}, alone and silent & Announce the
  homework and its form (packet or Desmos), circulate for the second formative read, preview
  Lesson~1.1 \\
\end{tabularx}
\end{fixedskillbox}

\begin{fixedskillbox}[Warm-Up --- Activate Prior Knowledge (5 min)]{frost}
\begin{tabularx}{\linewidth}{@{}X X@{}}
\begin{minipage}[t]{\linewidth}\vspace{0pt}
    \textbf{The three items and what each seeds}
    \begin{itemize}[leftmargin=1.1em, itemsep=1pt, topsep=2pt]
      \item \textbf{1.} $84$ of $240$ as a percent. \textbf{Seeds:} the summary a
            \emph{categorical} variable gets --- problem~6 (row~3, \emph{Two Types}) turns
            $3$ of $6$ captains into $50\%$ the same way.
      \item \textbf{2.} $18$ of $50$ scaled up to $900$. \textbf{Seeds:} sample $\to$
            population, which Lesson~1.2 names and homework item~5 spirals ($45$ of $300$,
            scaled to $2000$). Wordless about the vocabulary on purpose.
      \item \textbf{3.} A mean of five numbers, then a sentence naming the group and the
            units. \textbf{Seeds:} row~1 (\emph{Data in Context}) --- problem~1 asks for
            exactly this sentence again, and the ``$14$'' is the number the hook turns on.
    \end{itemize}
\end{minipage}
&
\begin{minipage}[t]{\linewidth}\vspace{0pt}
    \textbf{Running it}
    \begin{itemize}[leftmargin=1.1em, itemsep=1pt, topsep=2pt]
      \item \textbf{Debrief item~3(b) aloud, and only 3(b).} Take two student sentences and
            ask the class which one you could act on. Do not accept the bare ``$14$'' --- the
            sentence, with its group and its units, is what row~1 is built on. Write the best
            one on the board and leave it up all period.
      \item Item~3(b) is the tell. A student who writes ``$14$ minutes'' and stops has the
            units and not the who; row~1's display (A, B, C) is for that student.
      \item Item~2 separates the class: a student who adds instead of scaling the proportion
            needs a quiet minute during Guided Practice, not now. Items 1--2 should be quick
            --- say they are Algebra~1, so nobody thinks they are being tested.
    \end{itemize}
\end{minipage}
\end{tabularx}
\end{fixedskillbox}

\boxguard
\begin{skillbox}[Hook --- before anyone sits down]{frost}
The hook is on \textbf{page 1 of the notes}, under the vocabulary box, and on the board:
\textit{``The Riverbend roster's mean jersey number is $14$. That is low for a varsity squad
--- this is a young, inexperienced team.''} \textbf{Say:} every number in that sentence is
correct --- and let a student verify $84/6 = 14$ from the roster on the deck. \textbf{Then
ask: is the conclusion correct?} Students circle \emph{yes}, \emph{no}, or \emph{can't
tell} and write one line of reason in the hook box \emph{before} anyone speaks; then take the
hand vote and write the split on the board. Most rooms split between \emph{can't tell} and
\emph{yes}; a few say ``jersey numbers don't mean anything.'' Write that one down without
comment.
\textbf{Do not settle it.} Say only that the answer is problem~10, in the \emph{Digit Test}
row, and that they will vote again there. The vote is what makes the crux land: students
have to have accepted a correct computation as evidence before the notes take it away.
\end{skillbox}

\boxguard
\begin{skillbox}[Guided Notes \& Practice --- I do, then we do (35 min)]{frost}
\textbf{This is the whole lesson.} There is no group activity, and there is no independent
practice set on this page: \textbf{the homework is the individual practice}, started in class.
The notes are \textbf{one Main Ideas / Notes table}: four instruction rows (problems 1--13),
then the Guided Practice row (14--17). \textbf{The rhythm in every row is the same}: read the
definition aloud, mark up the display, work the first problem of the grid yourself, then hand
the rest to the room with the pen in their hand. The vocabulary box on page~1 is filled
\emph{as each term is named}, never in advance. The only blanks are the four cells of the
classify table (problem~5); every other answer is written in a problem's own space.
\begin{multicols}{2}
\textbf{I do / we do, row by row (about 20 min).}

\emph{Row 1, Data in Context (problems 1--2).} Write $12$ on the board and ask what it tells
them. Nothing. Read the definition, then have students name the three tags on the display ---
A is the \emph{who}, B the \emph{what}, C the \emph{units}. Work problem~1 yourself from the
warm-up sentence already on the board; the class does problem~2, which plants the crux: the
blog's $14$ has a who and a what and \emph{no units}, because a jersey number is not an
amount. Do not say more than that. Four minutes.

\emph{Row 2, Individual \& Variable (problems 3--4).} Put your finger on Nia's shaded row and
say ``one individual''; run it down the shaded Points column and say ``one variable.'' Work
problem~3 yourself and count aloud: $6$ and $5$ --- and address ``Player'' explicitly: it
names the rows, it is not a characteristic under study. The class does problem~4 (a column
and a player added: $7$ and $6$). Three minutes.

\emph{Row 3, Two Types (problems 5--9).} Read the two definitions and name the two panels,
then the two label boxes (Height is quantitative, Captain is categorical). Fill the classify
table (problem~5) with the room --- Position and Points first, opposite poles, then Captain
and Height. Work problem~6 yourself ($50\%$ captains --- a categorical variable does get
summarized, just not with a mean); the class does 7 and~8. \textbf{Take problem~9 last, and
take a hand vote before anyone writes}: Jersey \#, categorical or quantitative? Record the
split on the board and \textbf{do not settle it}. Seven minutes.

\emph{Row 4, The Digit Test (problems 10--13) --- the crux.} Put the number line up: six
players, a mean at $14$, nobody standing there. Ask \textbf{``find me the player who wears
$14$ on average''} and wait for the silence. Read the printed test, have students restate it
in the \emph{In your own words} space, and then \textbf{re-take the hook vote at problem~10
before you say anything}. Work problem~11 yourself (grade $10.5$); the class does 12 and~13.
Problem~13 is the counterweight --- Height is digits and averages fine. Six minutes.

\columnbreak
\textbf{We do (about 15 min).} The Guided Practice row, \emph{Student Life Survey} (problems
14--17) --- the same moves the homework will ask alone, on a survey instead of a roster.
\textbf{Students hold the pen; you hold the questions.} Problem~14 goes on them cold (count,
and say what you counted). \textbf{Problem~15 is the crux transferred}: Grade is $10, 11, 9,
12, 10, 11$, Theo says ``it's numbers,'' and $63/6 = 10.5$ puts nobody in grade $10.5$. Take a
hand count on Grade before anyone writes, and make a student --- not you --- compute the mean
aloud. Problem~16 is warm-up arithmetic ($108/6 = 18$); the sentence is the point. Problem~17
is the statistical question (PS.DC.1b): which one anticipates variability, and a rewrite.

Circulate during Guided Practice and demand the reason, not the word:
\begin{itemize}[leftmargin=1.1em, nosep]
  \item \textbf{Problem 14:} ``Point at one individual. Now point at one variable. Is
        \emph{Student} a variable?''
  \item \textbf{Problem 15:} ``Say it as a sentence --- \emph{Marisol is in grade 9.} Is 9 an
        amount of something? Then what would $10.5$ describe?''
  \item \textbf{Problem 16:} ``Who does the $18$ describe? What is it $18$ of?''
  \item \textbf{Problem 17:} ``How many answers does your question have? One look-up, or
        many?''
\end{itemize}
While circulating, pick the papers to put up in the debrief --- one problem~15 that fixed
Theo's reason with the test, and one that still says ``it's numbers.'' \textbf{Your formative
read comes twice}: here, and again in the last ten minutes as students start the homework.
\end{multicols}
\end{skillbox}

\boxguard[24]
\begin{skillbox}[Debrief --- whole class, spoken (10 min)]{frost}
Pens down, papers up. \textbf{This debrief is spoken} --- there is no box at the end of the
notes to fill in, so it lives entirely on the board and in the room.
\begin{multicols}{2}
\begin{enumerate}[leftmargin=1.3em, nosep, itemsep=3pt]
  \item \textbf{Put the blog's $14$ back up} and make a student say the sentence: the
        arithmetic is right, the variable is a label, the mean describes nobody. Then state it
        as the rule: \emph{a correct computation on the wrong kind of variable is still a
        worthless number.}
  \item Put up the two problem~15s you picked. The difference between them is not the
        verdict --- both say categorical --- it is the \emph{reason}. Read both aloud and ask
        which one would also catch Member ID tonight.
  \item Put problem~13 back up and take the over-correction: \emph{``Height is digits too.
        Who does $70$ inches describe?''} Digits are not the enemy.
  \item Take the statistical question last, from problem~17: how many answers does
        \emph{How long is Priya's commute?} have? One. That is why it takes no data.
\end{enumerate}
\columnbreak
\textbf{Check for understanding --- ask it cold, whole class.} \emph{``At Saturday's
cross-country meet, five Riverbend runners wore bibs $31, 18, 45, 9, 12$. A recap blog wrote:
`Riverbend's mean bib number was $23$ --- a low number for a varsity meet. This is a young,
inexperienced squad.' The arithmetic is right. Is the conclusion?''} Hands down, thirty
seconds of think time, then cold-call three students. Correct answer: \textbf{no} --- Bib \#
is categorical, a label written with digits; $115/5 = 23$ is correct and describes no runner
and no age. \emph{This replaces the exit ticket.}

\textbf{Formative read.} Sort the three answers as they come: \textbf{(a)} said \emph{no} and
named the type and the test; \textbf{(b)} said \emph{no} but reasoned ``the blog is biased''
or ``$23$ isn't low''; \textbf{(c)} said \emph{yes} or ``the math is right, so\ldots'' Pile
(b) is the teachable one and will be large --- they distrust the sentence without knowing
why; give them the word \emph{label}. \textbf{If pile (c) runs past a third of the room, open
Lesson~1.1 by re-running the \emph{Digit Test} row (problems 10--13)}, and say so plainly.

\textbf{If the debrief runs short}, cut items 3 and~4 --- item~1 is the one that cannot be
cut. \textbf{Do not let the debrief eat the last ten minutes} --- the homework start is not
optional padding, it is your second formative read.
\end{multicols}
\end{skillbox}

\boxguard
\begin{skillbox}[Homework --- scored, started in class, due after two study halls]{goldbox}
Six exercises on the \textbf{Northgate Recreation Center} ($8$ members on a record card):
\textbf{1} count the individuals and variables (\emph{the core procedure}); \textbf{2}
classify four columns with a reason --- \emph{Member ID} against \emph{Age} is the
\textbf{contrast pair}, two columns of digits, one a label and one an amount; \textbf{3} the
mean visits ($48/8 = 6$) written as data (\emph{interpret in context}); \textbf{4} the
percent taking a group class ($5/8 = 62.5\%$) and why a percent, not a mean
(\emph{justification}); \textbf{5} the \textbf{spiral} to Algebra~1 --- $45$ of $300$ as a
percent, scaled to $2000$, with a statistical question framing it; \textbf{6} the mean Member
ID of $140.5$ --- \emph{the crux head-on}, and it cannot be answered by computing anything.
The homework is capped at \textbf{two pages} (user decision 2026-09-06).

\textbf{Start it in the last ten minutes of the period, in class and alone.} \textbf{Scoring:}
12 points --- 2 per item. \textbf{Item~6 is where the misconception shows}: a student who
writes anything \emph{about} $140.5$ --- ``the typical member,'' ``the middle ID'' --- has
learned the word and not the rule. \textbf{Item~2's Member ID row carries the second check.}

\textbf{Packet or Desmos --- decided here.} The packet pages are authored and are the default.
Desmos Classroom has card sorts for categorical versus quantitative, but nothing that asks a
student to explain why a correctly computed mean measures nothing --- the item this lesson
exists for --- so \textbf{1.0 is a packet night}. Say so aloud at Close \& Assign.

\textbf{Preview of Lesson~1.1.} We name the four stages of the data cycle formally, write
questions a data set can actually answer, and organize collected answers into a frequency
table. Today's words --- individual, variable, categorical, quantitative --- are assumed from
here on.

\textbf{Connections \& Big Ideas.} \textit{Unit 1 core idea:} a conclusion is only as
trustworthy as the study that produced it. \textit{Standards covered:} PS.DC.1a, PS.DC.1b,
PS.DC.1c; AFDA.DA.2a. \textit{Spirals forward to:} population vs.\ sample and parameter vs.\
statistic (1.2, PS.DC.1d--e), sampling technique (1.3, PS.DC.2), bias (1.4), and
observational studies vs.\ experiments (1.5--1.7, PS.DC.3). \textit{Reaches back to:}
Algebra~1 percent and proportional reasoning, and the mean.
\end{skillbox}

\boxguard
\begin{skillbox}[Watch For (while circulating)]{redbox}
\begin{itemize}[leftmargin=1.2em, nosep]
  \item \textbf{The digit reflex} (problem~9, problem~15, HW~2 Member ID row, HW~6 --- the
        target misconception). Probe: \emph{``Compute the mean. Now find me the individual it
        describes.''}
  \item \textbf{Reasons that are not reasons} --- ``it's words,'' ``it's numbers'' (problem~5,
        problem~15, HW~2). Reject them all period; they are exactly the reasoning the digit
        test replaces. Probe: \emph{``Is one entry an amount of something, or a name for
        something?''}
  \item \textbf{Over-correction after the test} (problem~13, HW~2 Age and Visits rows): once
        Jersey \# is settled, Height and Commute get reclassified as categorical too. Probe:
        \emph{``Is $70$ inches an amount of anything?''}
  \item \textbf{Context-free answers} (problems~1, 7, 16, HW~3): a student who writes ``$18$''
        and stops has not finished. Probe: \emph{``Who does that describe, and $18$ of
        what?''}
  \item \textbf{The name column counted as a variable} (problems~3 and~14, HW~1): ``$6$
        variables.'' Probe: \emph{``Is \emph{Player} something we are studying about each
        person, or the way we tell the rows apart?''}
  \item \textbf{A one-answer question called statistical} (problem~17, HW~5): ``How long is
        Priya's commute?'' Probe: \emph{``How many answers does it have?''}
\end{itemize}
Cold-call prompts: \emph{``Name one individual in this table.''} \quad \emph{``Give me a
column of digits that is \emph{not} quantitative --- and one that is.''} \quad \emph{``Change
this question so it takes data to answer.''}
\end{skillbox}

\boxguard
\begin{skillbox}[Close \& Assign (10 min)]{goldbox}
\textbf{Students start the homework here, in class, alone and silent --- this is the individual
practice, and the first minutes of it are your best formative read.} Hand the launch first ---
six exercises on paper, scored out of 12, due at the start of the first class after two study
halls --- and point out that the \emph{Keep in Mind} box on the cover is the page to flip back
to. Circulate: \textbf{item~2's Member ID row and item~6 are the lines that tell you whether
the \emph{Digit Test} row landed}, and item~6 is on page~2, so put students on 1, 2, and~6
first and let 3, 4, and~5 go home. Sort what you see into three piles as you walk --- called
Member ID categorical \emph{and} gave the amount-versus-label reason (ready); called it
categorical with ``it's an ID'' and no reason (ready, needs the language); called it
quantitative or wrote something about $140.5$ (the misconception survived) --- and open
Lesson~1.1 by re-running problems 10--13 if that third pile ran past a third of the room. Then
close on what changed: \emph{``You walked in able to compute a mean. You walk out able to say
when a mean is allowed to exist --- and that is decided by the variable, not by the
arithmetic.''} \textbf{Preview:} Lesson~1.1, \emph{The Data Cycle} --- today's words are
assumed from here on; next time we write the questions a data set can actually answer.
\end{skillbox}

% ── Teacher notes — one per component, in packet order (three) ──────────────
% Teacher-only prose lives HERE, never in a _key: a note is the one block in a key with no
% counterpart in the blank. No note for the debrief or the homework start — both are fully
% specified in their own boxes above.
\begin{teachernote}[Warm-Up]
Five minutes, silent, timed. All three items are Algebra~1 review and should be quick; the
point is diagnostic, not instructional. \textbf{Item~3(b) is the one to read while
collecting}: students who write a sentence with a group and units have already made row~1's
distinction, and row~1 can move fast. Students who write ``$14$'' or ``$14$ minutes'' and
stop are the ones the A/B/C display is for. Write the best sentence on the board in a
student's own words and \textbf{leave it up all period}; problem~1 asks for it again, and the
debrief circles back to it. Item~2 (scaling $18/50$ up to $900$) is the one that separates
the class --- a student who adds or multiplies $18 \times 18$ needs a quiet minute during
Guided Practice. Do \emph{not} supply the words \emph{sample} and \emph{population} here ---
Lesson~1.2 names them.
\end{teachernote}

\begin{teachernote}[Guided Notes \& Practice]
\textbf{This one document is the lesson --- pace it 20 / 15, then 10 for the debrief, then 10
for the homework start.} It is one table, and every row runs the same way: you read the
definition, mark up the display, and work the first problem of the grid (1, 3, 6, 11); the
class works the rest with the pen in their hand. Rows 1 and~2 must move: four minutes on data
in context, three on individuals and variables. Problem~2 plants the crux --- the blog's $14$
has no units --- and you should let it sit unexplained.

\textbf{The single highest-leverage move today is in row~3}: classify Jersey \# \emph{last}
(problem~9), take a hand vote, record the split on the board, and then \textbf{do not settle
it}. The digit test in row~4 lands because students have already committed to a side and are
about to watch a correct computation produce a worthless number. If you resolve Jersey \#
during row~3, row~4 becomes a definition to copy rather than an argument to win, and the
misconception survives the lesson.

\textbf{Spend the time in row~4, \emph{The Digit Test}}, which carries the misconception. Put
the number line up, ask for the player who wears $14$ on average, and wait for the silence.
Then \textbf{re-take the hook vote at problem~10 before you say anything}: yes or no. Let the
room argue. Do not resolve it by repeating the definition; count the players on the number
line and point at $14$ --- \emph{nobody}. Problem~13 is the counterweight: Height is digits
and averages fine, and a room that has just learned the test will try to fail Height with it.

\textbf{Guided Practice (14--17) is where the pen changes hands.} \emph{Student Life Survey}
runs the four moves the homework will ask alone --- count and name, classify with the test,
compute and say it as data, recognize a statistical question --- on a survey instead of a
roster. Problem~15 is the crux transferred: Grade is $10, 11, 9, 12, 10, 11$, and Theo's
``it's numbers'' is the coach's reason from the old meet-card activity, now on a student's
lips. Make a student compute $63/6 = 10.5$ aloud and ask who is in grade $10.5$. If you only
get through 14, 15, and~16, that is the right trade; 17 can be posed aloud.

\textbf{There is no independent practice set on this page, and that is deliberate --- the
homework is the individual practice.} Guided Practice is the last row; the period ends with
students starting the homework in front of you, which is where your second formative read
comes from. Item~2's Member ID row is the one to watch. If the ``it's numbers'' pile is more
than a third of the room, \textbf{open Lesson~1.1 by re-running problems 10--13}, and say so
plainly.

\textbf{Debrief (10 min), spoken.} The blog's $14$ back on the board with a student saying
\emph{correct arithmetic, wrong kind of variable, describes nobody}; it is the one move that
cannot be cut. The cold check is the meet-card bib claim --- the same trap on fresh digits.
\textbf{Do not let the debrief eat the last ten minutes} --- the homework start is not
optional padding.
\end{teachernote}

\begin{teachernote}[Homework]
\textbf{Scored, 12 points (2 per item), due at the start of the first class after two study
halls.} The ten minutes at the end of the period are a \emph{start}, not a completion: put
students on 1, 2, and~6 while you can still catch a wrong start, and let 3, 4, and~5 go home.
Item~6 is the one to read first when scoring --- any answer that computes or interprets
something \emph{about} $140.5$ is a miss, however well written. Item~2 is the fullest
diagnostic in the set: Member ID and Age are both digits, and a student who labels Member ID
quantitative, or labels Age categorical, is still sorting by what the column looks like. Score
items 1--5 for correctness and item~6 for the \emph{reason}, not the verdict. \textbf{When
you score the page, sort item~6 into three piles} --- said \emph{nothing} and named the type
(ready); said \emph{nothing} vaguely (ready, needs the language); wrote something about
$140.5$ (the misconception survived). If more than a third land in the last pile, reopen the
jersey number line in the Lesson~1.1 warm-up debrief rather than letting it ride.

\textbf{Packet is the call for 1.0.} Desmos carries the classification card sort and nothing
else in this set; do not swap it in. Go over item~2's Member ID row and item~6 at the start
of Lesson~1.1, in that order.
\end{teachernote}
\end{document}
